\documentclass[a4paper]{article}
\usepackage{amsmath}
\usepackage{amsfonts}

\usepackage{amssymb}
\usepackage{graphicx}
\usepackage{float}

\begin{document}

\noindent \large Auteur: Siebe Haché \\
\noindent \large Studierichting: Informatica \\
\noindent \large Oefeningenreeks 6A nr. 10.5 \\

\medskip
\normalsize

Bepaal een meetkudnige rij voor 5 termen $a, b, c, d$ en $e$, waarvoor $a - e = 60$ en $b - d = -24$.\\

$\left\{ \begin{array}{l} 
    a - ar^4 = 60 \Rightarrow a(1 - r^4) = 60\\

    ar - ar^3 = -24 \Rightarrow ar(1 - r^2) = -24\\
\end{array} \right. $ \\

eerste vergelijking delen door de tweede geeft: \\

$\dfrac{a(1 - r^4)}{ar(1 - r^2)} = \dfrac{60}{-24}$ \\

$\Rightarrow\dfrac{a(1 - r^2)(1 + r^2)}{ar(1 - r^2)} = -\dfrac{5}{2}$\\

$\Rightarrow\dfrac{1 + r^2}{r} = -\dfrac{5}{2}$ \\

$\Rightarrow 2(1 + r^2) = -5r$ \\

$\Rightarrow 2r^2 + 5r + 2 = 0$ \\

$D = 5^2 - 4 \cdot 2 \cdot 2 = 25 - 16 = 9$ \\

$r_{1} = \dfrac{-5 + 3}{4} = -\dfrac{2}{4} = -\dfrac{1}{2}$ \\

$r_{2} = \dfrac{-5 - 3}{4} = -\dfrac{8}{4} = -2$ \\

Geval 1: $r = -2$ \\

$a(-2)(1 - (-2)^2) = -24$ \\

$a(-2)(1 - 4) = -24$ \\

$a(-2)(-3) = -24$ \\

$6a = -24 \implies a = -4$ \\

Geval 2: $r = -\dfrac{1}{2}$ \\

$a(-\dfrac{1}{2})(1 - (-\dfrac{1}{2})^2) = -24$ \\

$a(-\dfrac{1}{2})(\dfrac{3}{4}) = -24$ \\

$-\dfrac{3}{8}a = -24 \implies a = \dfrac{-24 \cdot 8}{-3} = 64$ \\

\begin{tabular}{|c|c|c|c|c|c|c|} 
\hline 
Oplossing & $a$  & $b$   & $c$   & $d$  & $e$ \\
\hline 
1         & $-4$ & $8$   & $-16$ & $32$ & $-64$ \\
\hline 
2         & $64$ & $-32$ & $16$  & $-8$ & $4$ \\
\hline
\end{tabular} \\

\begin{thebibliography}{999}
\end{thebibliography}
\end{document}